\textcolor{red}{INSTRUÇÕES GERAIS:}

\textcolor{red}{Observação 1: Este documento é uma recomendação de formatação da monografia, contendo inclusive itens que são esperados no texto da monografia. Contudo, o aluno tem a liberdade de discutir com seu orientador e inserir outros capítulos/seções ou reestruturá-la, caso seja necessário.}

\textcolor{red}{Observação 2: O texto da monografia (marcado como estilo normal do Word)  deve utilizar fonte Times New Roman tamanho 12pt com espaçamento 1,5 entre linhas, conforme já configurado neste documento. Utilizar papel A4.}

\textcolor{red}{Observação 3: Margens: Superior e esquerda: 3,0 cm; Inferior e direita: 2,0 cm.}

\textcolor{red}{Observação 4: Sugere-se que a monografia tenha em torno de 30  páginas.}

\textcolor{red}{Observação 5: Outras recomendações: (i) Substituir todos os textos marcados em vermelho; e (ii) Seções que forem opcionais e não forem utilizadas na monografia devem ser retiradas.}

\textcolor{red}{Lista de Abreviaturas/Siglas: Esta seção é opcional. Insira-a na monografia caso haja uma quantidade razoável de abreviaturas. O fato de incluir a abreviatura/sigla aqui não o libera da obrigatoriedade de colocar a sigla e seu significado na primeira ocorrência dela no texto. A partir da segunda ocorrência use somente a sigla. Se o leitor se esquecer do significado, pode recorrer a esta lista pra se lembrar, ao invés de ter que ficar procurando no meio do texto.}

\textcolor{red}{Lista de Símbolos: Esta seção é opcional. Insira-a na monografia caso haja uma quantidade razoável de símbolos.}

\textcolor{red}{Lista de Gráficos: Esta seção é opcional. Insira-a na monografia caso haja gráficos ao decorrer do texto.
Sugestão: Utilize o índice de ilustrações (gráficos) do Word para gerar automaticamente esta lista.}

\textcolor{red}{Lista de Tabelas: Esta seção é opcional. Insira-a na monografia caso haja tabelas ao decorrer do texto.
Sugestão: Utilize o índice de ilustrações (tabelas) do Word para gerar automaticamente esta lista.}

\textcolor{red}{Lista de Algoritmos: Este item é opcional. Insira-o na monografia caso haja uma quantidade razoável de algoritmos.}

\textcolor{red}{Lista de Códigos-fonte: Este item é opcional. Insira-o na monografia caso haja uma quantidade razoável de códigos-fonte.}

\textcolor{red}{Lista de Figuras: Esta seção é opcional. Insira-a na monografia caso haja figuras ao decorrer do texto.
Sugestão: Utilize o índice de ilustrações (figuras) do Word para gerar automaticamente esta lista.}

\textcolor{red}{INSTRUÇÕES INTRODUÇÃO:}

\textcolor{red}{Observação: um bom capítulo de introdução tem em torno de 3 a 5 páginas. Sugere-se que a monografia tenha em torno de 30 páginas. Caso tenha muito conteúdo, pode escolher algumas partes e colocar como Apêndice. Substitua todos os textos marcados em vermelho e também as Seções que forem opcionais e não forem utilizadas.}

\section{Motivação e Contextualização}

\textcolor{red}{Nesta seção descrevem-se a área de pesquisa na qual o trabalho está inserido, o problema e/ou as circunstâncias que motivaram o projeto e as potenciais contribuições oriundas de sua realização. Além disso, é necessário sintetizar o que foi feito no Projeto I, caso o aluno esteja matriculado atualmente no Projeto II. Ao longo de todo o texto, se citar figuras, cite-a pelo nome, explique a figura e cite a fonte (se for sua, coloque Fonte: autor desta monografia), como nesse exemplo (veja a seção Corpos Flutuantes Figura \ref{figura:logomarca_usp}). Vejam que a legenda vai acima da Figura e a fonte abaixo. Sempre que se referir à figura, use seu número (e não: figura abaixo, figura acima, figura a seguir, etc). O mesmo vale para Tabela, Gráfico, Algoritmo, etc. Toda figura, tabela, etc. tem que estar citada no texto para chamar a atenção do leitor para ela.}

\section{Objetivos}

\textcolor{red}{Indique claramente quais são os objetivos do trabalho, caracterizando de forma sucinta o que se pretende atingir. Se o emprego de uma metodologia específica de trabalho for relevante, indique também; caso contrário, ela aparece apenas na seção de desenvolvimento do trabalho.}

\section{Organização}

\textcolor{red}{Escreva um parágrafo que resume cada um dos demais capítulos da monografia, fazendo referência a cada um deles – como no exemplo que segue. Indique também a existência de apêndices e anexos, se houver.  No Capítulo 2 é apresentada uma revisão da terminologia básica utilizada no projeto, bem como os trabalhos da literatura relacionados ao presente projeto. A seguir, no Capítulo 3, descrevem-se as atividades realizadas. Finalmente, no Capítulo X, apresentam-se as conclusões e trabalhos futuros. Observe que o tempo verbal é presente, pois refere-se ao texto que já está escrito na monografia. Use tempo futuro somente para se referir a atividades que ainda serão feitas.}